% ===== 原文 p.35 =====
% tr-05 片段 c：算子公式、紧致情形的 WP 体积例、讲稿收尾（p.35）
% 终点：写到 "…on the moduli space [HM98]." 为止；不写 References 标题与条目（归 tr-06）。
% 本块无编号公式；无编号陈述；无图。
% [原文 p.35]
$\pi:\Mgnb\longrightarrow\Mgnb$，即遗忘最后一个穿刺点的态射，则 $\Mgnb$ 上的类
$\wt{\kappa}_m,\wt{\psi}_k$ 与 $\Mgnb$ 上的 $\kappa_m,\psi_k$ 满足
$\wt{\kappa}_m=\pi^*\kappa_m+\psi_{n+1}^m$ 以及 $\wt{\psi}_k=\pi^*\psi_k,+D_k$
（其中 $D_k$ 是 $\Mgnb$ 上万有曲线里第 $k$ 个穿刺点的除子），[HM98]。

% [Tutor Note: p.35 的 $\wt{\psi}_k=\pi^*\psi_k,+D_k$ 在原文里 $\psi_k$ 与 $+D_k$ 之间确有一个多余逗号（已用 700 dpi 裁图复核，是原文如此，不是文字层识别错误），按上下文应读作 $\wt{\psi}_k=\pi^*\psi_k+D_k$；正文照原文保留。另外注意下标与上标之别：$\wt{\kappa}_m$、$\wt{\psi}_k$ 的下标 $m,k$ 是类的编号，而 $\psi_{n+1}^m$ 的 $m$ 是幂次——两者次数同为 $2m$，故 $\wt{\kappa}_m=\pi^*\kappa_m+\psi_{n+1}^m$ 是齐次等式。]

% [Tutor Note: 紧化模空间的上划线在原文里并不处处一致，容易抄错：p.31 的定义式、
%  p.34 的两个积分、p.35 的本段（$\pi:\Mgnb\to\Mgnb$、$D_k$ 所在的 $\Mgnb$）都带 overline；
%  而 p.31 正文的 "moduli spaces $\cM_{g,n}$" 与 p.32 的 $\int_{\cM_{1,1}}$ 确实**不带**。
%  本译本逐处按页图保留，未做"看起来更整齐"的统一。]

第二个证明基于体积递推中诸运算的精确公式。对广义 dilaton 方程，考虑作用在环
$\C[x^2,y^2,L_1^2,\dots,L_n^2,\dots]$ 上的下列四个算子
\[
2\frac{\partial}{\partial L_1}L_1[\cdot],\qquad
\frac{\partial}{\partial L_{n+1}}[\cdot]\Big|_{L_{n+1}=2\pi i},\qquad
\int_0^{\infty}\!\!\int_0^{\infty}xyH(x+y,L_1)[\cdot]\,\dd x\,\dd y,
\]
以及
\[
\int_0^{\infty}x\big(H(x,L_1+L_k)+H(x,L_1-L_k)\big)[\cdot]\,\dd x.
\]

% [Tutor Note: p.35 四个算子里的 $H$ 就是 p.18 公式 (1) 所定义的 $H(x,y)=\frac{1}{1+e^{(x+y)/2}}+\frac{1}{1+e^{(x-y)/2}}$，故后两个积分算子正是把体积递推中对 $H$ 的积分按 $L_1+y$ 与 $L_1\pm L_k$ 的变量替换展开。第一个算子 $2\frac{\partial}{\partial L_1}L_1[\cdot]$ 读作“先乘以 $L_1$ 再对 $L_1$ 求导”（算子从右向左作用），不要误读成先求导再乘 $L_1$：按乘积法则二者相差 $2[\cdot]$，即 $2\frac{\partial}{\partial L_1}(L_1 f)-2L_1\frac{\partial f}{\partial L_1}=2f$，这一项恰会改变后面的项合并。]

诸算子的公式由此展开。证明时，广义 dilaton 方程用第二个算子写出，对左端应用体积
递推，再应用算子公式，合并各项即得右端。广义 dilaton 方程给出紧致情形的 WP 体积，
包括如下例子
\[
V_{2,0}=\frac{43\pi^6}{2160},\qquad
V_{3,0}=\frac{176557\pi^{12}}{1209600},
\]
\[
V_{4,0}=\frac{1959225867017\pi^{18}}{493807104000}
\qquad\text{以及}\qquad
V_{5,0}=\frac{84374265930915479\pi^{24}}{355541114880000}.
\]

% [Tutor Note: p.35 四个体积值可作量纲自检：$V_{g,0}$ 的总次数为 $6g-6+2n=6g-6$，故 $g=2,3,4,5$ 应分别带 $\pi^6,\pi^{12},\pi^{18},\pi^{24}$，与原文逐项相符（这一列指数是最容易抄错的地方）。四个分式的分子分母已用 600–700 dpi 裁图逐位核对：$43/2160$、$176557/1209600$、$1959225867017/493807104000$、$84374265930915479/355541114880000$，与文字层一致。]

我们以如下观察结束本讲：关于模空间上的 Witten 猜想，研究十分广泛。Kontsevich 与
Mirzakhani 的途径本质上是分析的。Okounkov 与 Pandharipande [OP09] 把问题转化为计数
Hurwitz 数——即球面分支覆盖的拓扑型——并用组合方法计数置换分解为对换的分解。他们的
组合途径给出 Kontsevich 的公式。关于模空间上的相交数与关系，还有更大规模的研究 [HM98]。
